% Kata pengantar
\begin{center}
{\bfseries\large KATA PENGANTAR}
\end{center}

\vspace{0.5cm}

Laporan Tugas Akhir berjudul \textit{Pengembangan Companionship Chatbot dengan Asesmen Depresi melalui Percakapan dan Memori Jangka Panjang Berbasis Knowledge Graph untuk Mahasiswa Indonesia} ini disusun sebagai salah satu syarat untuk memperoleh gelar Sarjana pada Program Studi Teknik Informatika, Sekolah Teknik Elektro dan Informatika, Institut Teknologi Bandung, sekaligus sebagai luaran akhir mata kuliah IF4092 Tugas Akhir II.

Selama proses penyusunan tugas akhir ini, penulis memperoleh banyak bantuan, bimbingan, dukungan, serta pengalaman berharga dari berbagai pihak. Oleh karena itu, penulis menyampaikan terima kasih kepada:

\begin{enumerate}[label=\arabic*.,leftmargin=*]

\item Dr. Agung Dewandaru, S.T., M.Sc. selaku dosen pembimbing yang telah meluangkan waktu dan tenaga untuk memberikan arahan, masukan, dan berbagai diskusi ilmiah selama proses penyusunan tugas akhir ini.

\item Seluruh dosen dan tenaga kependidikan Program Studi Teknik Informatika ITB yang telah memberikan ilmu, pengalaman, serta bantuan administrasi selama masa perkuliahan.

\item Seluruh dosen serta rekan rekan asisten di Laboratorium Ilmu Rekayasa dan Komputasi, Teknik Informatika ITB, atas kesempatan, pengalaman, diskusi, serta lingkungan belajar yang sangat berharga selama penulis menjalankan tugas sebagai asisten laboratorium.

\item Kedua orang tua serta seluruh keluarga penulis yang senantiasa memberikan dukungan, perhatian, motivasi, dan semangat selama masa perkuliahan hingga penyelesaian tugas akhir ini.

\item Teman teman seperjuangan di Program Studi Teknik Informatika ITB yang telah berbagi ilmu, pengalaman, dan kebersamaan selama menjalani perkuliahan.

\item Teman teman kos Dago Asri A10 yang telah menjadi tempat berbagi cerita, saling membantu, dan memberikan dukungan selama menjalani kehidupan perkuliahan.

\item Teman teman SMA 3 Bandung yang senantiasa memberikan semangat dan menjaga tali silaturahmi meskipun menempuh perjalanan yang berbeda.

\item Seluruh pihak lain yang tidak dapat penulis sebutkan satu per satu, yang telah memberikan bantuan maupun dukungan dalam penyelesaian tugas akhir ini.

\end{enumerate}

Penulis menyadari bahwa laporan ini masih memiliki berbagai keterbatasan. Oleh karena itu, kritik dan saran yang membangun sangat diharapkan sebagai bahan evaluasi dan pengembangan pada penelitian selanjutnya. Penulis berharap laporan ini dapat memberikan manfaat bagi pembaca serta menjadi salah satu referensi bagi pengembangan penelitian di bidang rekayasa perangkat lunak, kecerdasan buatan, dan sistem berbasis knowledge graph.

\vspace{1.5cm}

\begin{flushright}
Bandung, \rule{4cm}{0.4pt}\\[1cm]
Penulis
\end{flushright}
